% CREATED BY DAVID FRISK, 2016
\chapter{Background}
\label{chap:background}

% 2.0 Short Introduction to Theory Chapters - Outline the structure of these chapters and the purpose for each. 
This section introduces necessary context for further discussion within this thesis. The in-vehicle system and related components will be introduced, followed by an introduction to modern challenges in regard to data processing and data usage in the automotive industry. The purpose of this chapter is to provide the necessary background for the reader to understand the motivation and relevance of the research presented in this thesis. It also serves to establish a common understanding of key concepts and terminology that will be used throughout the thesis.

% 3.1 In-Vehicle Systems and Event-Triggered Logging - DONE (TIM) - DONE (EMRIK)
\section{In-Vehicle Systems and Event-Triggered Logging}

\textbf{The In-Vehicle Embedded System} An embedded system is a computer system integrated into a physical device or product to perform a dedicated function as part of a larger system~\cite[pp.~vii--x,~1--3]{Lee2017}. The in-vehicle embedded system is a specialized embedded system integrated within a vehicle. It is essential for real-time functions like controlling, monitoring, and enhancing various automotive operations~\cite[pp.~1-1--1-5]{Navet2009}. These in-vehicle embedded systems, like embedded systems in general, typically consist of both hardware and software components, such as electronic control units (ECUs), sensors, actuators, and communication interfaces~\cite[pp.~3--4]{Lee2017},~\cite[pp.~1-3--1-10]{Navet2009}. 

ECUs serve as the primary processing nodes in in-vehicle embedded systems. They are responsible for executing control algorithms on sensor inputs and issuing commands to actuators. They also handle local control logic while coordinating via in-vehicle networks to form a distributed system~\cite[pp.~1-8--1-19]{Navet2009}. An example of how ECUs work together to create the in-vehicle embedded system is illustrated by~\cite[pp.~1-8--1-9]{Navet2009}. In the example an ECU coordinates a request, e.g., locking or unlocking a door, while the controller realize the request on the physical system, which is controlled by another ECU\@. The resulting system is the main generator of telemetry data within the vehicle. Vehicle telemetry data, usually defined as data that is produced by the vehicle and sent off-board, is key to improving important aspects of these vehicles like safety~\cite[pp.~1--5]{MITRE2023}. This vehicle telemetry data is at the core of this thesis.

\textbf{In-Vehicle Networks} As mentioned above, in-vehicle networks serve as the connectors between the ECUs of the in-vehicle embedded system. 
Among these in-vehicle networks Control Area Networks (CAN), Local Interconnect Networks (LIN), FlexRay, and Ethernet are popular representatives that enable efficient communication and coordination among different vehicle subsystems~\cite[pp.~1-2--1-5]{Navet2009}. 

From a technical perspective, ECUs connect to these networks via dedicated communication services. The communication services are tasked with providing a uniform interface to the network, while hiding protocol and message properties from the application~\cite[p.~48]{AUTOSAR2022}. The in-vehicle network protocols are categorized into three classes, A, B, and C. LIN falls into category A with a bit rate below 10 kbps and application in sensor and actuator networks. CAN-B is a representative of category B with a bit rate between 10--500 kbps. Category B networks serve vehicle-centric domains and body electronics. Category C in-vehicle networks like FlexRay or CAN-C are defined by having a bit rate of lower than 1 Mbps and used in safety-relevant systems~\cite[pp.~1-12--1-13]{Navet2009}. Despite their effectiveness in their respective use-cases, all in-vehicle networks must contend with fundamental bandwidth limitations that constrain data throughput across the vehicle~\cite[pp.~1-18--1-19]{Navet2009}.

\textbf{Event-Triggered and Time-Triggered Logging} Event-triggered logging and diagnostic frameworks, which record data only when specific error codes occur or threshold crossings are detected, are widely adopted to reduce data transmission and avoid bus saturation in complex systems such as the in-vehicle embedded system~\cite[pp.~7--13]{AUTOSAR2022b}. According to the AUTOSAR Diagnostic Log and Trace (DLT) standard, this mechanism can be implemented as follows: When software components or Basic Software modules report errors via the Default Error Tracer, the DLT module receives these via a specified function and stores timestamped log/trace messages to flash storage~\cite[p.~25]{AUTOSAR2022b}. Non-critical events are filtered during normal operation~\cite[pp.~30--35]{AUTOSAR2022b}. 

Time-triggered logging serves as an alternative to event-triggered logging and works by setting predetermined points in time where data is transmitted. The validity of this strategy is easily monitored, as the time-frame scheduling is predictable and missing messages are immediately identified~\cite[pp.~4-5]{Navet2009}. Although Navet discusses these paradigms in the context of communication scheduling, similar principles apply to diagnostic event- and time-based mechanisms, which leads to similar struggles in these systems.

While these approaches have been effective in reducing the data load of the in-vehicle networks, they come with major limitations. Two immediate issues these diagnostic frameworks struggle with is presenting an unbiased and comprehensive monitoring of the system, while having to carefully tune event thresholds and diagnostic parameters~\cite[pp.~4-5--4-6]{Navet2009},~\cite{EventDrivenDiagPatent}. In the next section we will go further into how these limitations are accentuated by recent developments. 

% 3.2 Downstream ML Tasks and Data Quantity - DONE (TIM) - DONE (EMRIK)
\section{Downstream Machine Learning Tasks and Data Quantity}\label{sec:downstream-ml-and-dq} 

Two developments in recent years further underline the shortcomings of event-triggered logging in automotive systems: the massive increase in signal-based data in the in-vehicle network and the growing relevance of downstream machine learning (ML) tasks. 

\textbf{Data Quantity in Modern Vehicles} Recent industry and research reports indicate that the data quantity generated by sensors and the amount of interconnectivity in vehicles is growing at an extremely rapid pace. According to a 2021 automotive electronics report, by 2030, about 95\% of new vehicles will be connected, up from around 50\% today~\cite{Bertoncello2021}. Another report states that even today, a single car can generate up to 1 terabyte (TB) of data per hour from its sensors~\cite{Samantaray2023}. This explosive growth is driven by the increasing number and sophistication of sensors—such as cameras, radars, and lidars—required for advanced safety and autonomous driving features, with the complexity and volume of data presenting significant challenges for storage, processing, and transmission within embedded automotive systems~\cite{Samantaray2023},~\cite[pp.~3--5]{Komorkiewicz2023}. While much of the focus in research and reporting on this issue centers on autonomous driving technologies like lidars and cameras, it is important to note that the core of the in-vehicle system, the ECU, is also a major contributor to the data quantity issue. This implies that not only autonomous driving technologies but modern vehicles in general are facing these challenges~\cite[pp.~3--4]{Komorkiewicz2023}.

\textbf{Downstream ML Tasks} Modern vehicles increasingly rely on data-driven intelligence to enhance safety, reliability, and efficiency. As there currently exist no conclusive definition of these tasks, we will refer to them as downstream ML tasks within the context of this thesis. These downstream ML tasks we define as the offline use of collected vehicle data for the training of ML models for various tasks. This is distinctly different from real-time perception and control systems. Downstream ML tasks have become central to automotive-system design. Theissler et al.~\cite[p.~2]{Theissler2021} mention three reasons for the rise of importance of these downstream ML tasks: data availability, breakthroughs in algorithm development and the advancements in computational power. Prominent examples of downstream ML tasks in the automotive context include: Predictive maintenance tasks aim to predict the expected time of a failure, thereby helping to estimate remaining functionality of components or systems~\cite[p.~3]{Theissler2021}. Anomaly and intrusion detection tasks want to find anomalies or outliers in data points which is especially relevant in fragile systems like the CAN network~\cite[pp.~1--3]{Oezdemir2024}. Fleet-level management operations like fuel consumption that aim to better utilize important assets in vehicle systems~\cite[pp.~1--2]{Brunheroto2022}. Utility loss in the datasets used to train the models for these tasks translates directly to degraded model performance, increased operational costs, and missed business opportunities~\cite[pp.~4--5]{Komorkiewicz2023}. This further underlines the importance of a comprehensive data logging strategy.

There exists relatively little research on the use of event-triggered logging systems as a compression strategy and it's impact on downstream ML tasks. It is therefore difficult to definitively say that event-triggered logging systems specifically produce data with lower utility. Studies have however shown that similar utility agnostic data reduction strategies can have negative impacts on downstream utility in the domain of time-series data. One such example comes from a study by~\cite{Tirupathi2022} where the authors show that a 10-minute aggregation of the data leads to degraded utility for downstream ML tasks. In this specific case, the F1-score of the predictive downstream model was reduced by from 0.63 to 0.50 when trained on the aggregated data compared to the original dataset. The same study also shows that a model based lossy compression technique can be implemented without significant utility loss which suggests that more strategic data reduction techniques can retain more utility for downstream ML tasks. Further analysis of the impact of lossy compression algorithms on time series task utility in general has been done by~\cite{Eichinger2015} and~\cite{Moon2018}. It is important to note, that, while challenging, at specific error bounds compression techniques within these studies can actually show improvement to the task metrics. All this highlights the need for a strategic lossy compression technique that accounts not only for reconstruction, but also for downstream task utility. 

\section{Related Work}

\textbf{Time-Series Neural Compression} Neural compression for time-series data has been studied, but many methods are not optimized for edge deployment. The architecture proposed by Zheng and Zhang~\cite{Zheng2023} is representative of this line of work. Their method uses a continuous-latent autoencoder with recurrent modules and is primarily optimized for reconstruction quality, while computational efficiency and downstream task utility retention are not central optimization targets~\cite{Zheng2023}.

\textbf{Edge-Optimized Neural Compression} The majority of research labeled as ``edge-optimized neural compression'' focuses on compressing ML models for deployment rather than compressing input data. Lai et al.~\cite{Lai2026} highlight the importance of enabling complex deep learning models on resource-constrained edge devices. Key developments include improved efficient transformer variants and knowledge-distillation methods~\cite{Lai2026}. These advances address one side of the edge-optimization problem (model compression), while the complementary side is efficient compression of the data itself. EdgeCodec is a notable example of this latter direction and is discussed below.

\textbf{Audio/Speech Neural Compression} The audio/speech community contains extensive research on neural compression, with a strong emphasis on vector quantization. A key development in this literature is residual vector quantization (RVQ). RVQ works by layering hierarchical quantizers whith each layer quantizing the residuals of the previous layer~\cite[pp.~495--507]{Zeghidour2021}.

\begin{itemize}
    \item SoundStream (Google, 2021): first major use of RVQ for audio compression~\cite[pp.~495--507]{Zeghidour2021}.
    \item EnCodec (Meta, 2022): refined RVQ for high-quality audio~\cite{Defossez2022}.
    \item WaveTokenizer (2024): recent RVQ-based audio codec~\cite{Ji2025}.
\end{itemize}

\textbf{EdgeCodec} EdgeCodec transfers RVQ ideas from audio/speech processing to time-series compression (aerospace sensor data) and is explicitly designed for edge deployment. This is achieved through a lightweight encoder deployed on the edge and a heavier decoder executed in the cloud~\cite{Hodo2025}. As a result, EdgeCodec is a notable outlier in neural compression literature and addresses many challenges apparent in the research field. However, it has important limitations for the present use case. For one, EdgeCodec is optimized for smooth, continuous aerospace sensor streams rather than mutlimodal vehicle telemetry. And second, EdgeCodec focuses on reconstruction error and omits analysis of downstream task utility retention.

The literature review indicates that neural compression, VQ-based quantization, and lightweight design principles have all been explored individually. The remaining gap is a method tailored to mutlimodal data that explicitly optimizes downstream ML task utility rather than reconstruction quality alone. Based on this gap, the present thesis evaluates two existing solutions, namely the continuous-latent autoencoder by Zheng and Zhang~\cite{Zheng2023} and EdgeCodec~\cite{Hodo2025}, with respect to reconstruction quality on multimodal data and downstream ML task performance. Baseline models derived from these works are implemented and evaluated on automotive and IoT telemetry data, with the objective of identifying and improving design choices for multimodal data and downstream utility retention.
